% mazzi: 03,04
\chapter{Norme sportive antidoping (NSA)}

% ---------------------------------------------------------------------------
\section{Il Codice Sportivo Antidoping}

\textbf{Codice Sportivo Antidoping} (\textbf{CSA}): attuativo del Codice WADA, disciplina la
definizione di doping e le fattispecie che ne costituiscono violazione.

% ---------------------------------------------------------------------------
\section{Articolo 1 -- Definizione di doping}

Per doping si intende la violazione di una o più norme contenute nei successivi articoli 2 e 3
del presente Codice.

% ---------------------------------------------------------------------------
\section{Articolo 2 -- Violazioni del Codice WADA}

Per doping si intende la violazione di una o più norme contenute negli articoli dal 2.1 al 2.10
del presente Codice. Lo scopo dell'articolo 2 è quello di specificare le circostanze e la
condotta che costituiscono violazione delle norme antidoping: i procedimenti nei casi di doping
saranno basati sul presupposto che una o più di queste specifiche norme siano state violate. Gli
\textbf{Atleti} o altre Persone saranno tenuti a conoscere cosa costituisca violazione delle
norme antidoping e le sostanze vietate e i metodi proibiti inclusi nella Lista.

\subsection{2.1 -- La presenza di una sostanza vietata}

\textbf{Art.\ 2.1}: la presenza di una sostanza vietata o dei suoi metaboliti o marker nel
campione biologico dell'Atleta.
\begin{itemize}
  \item \textbf{2.1.1}: ciascun Atleta deve accertarsi personalmente di non assumere alcuna
    sostanza vietata, poiché sarà ritenuto responsabile per il solo rinvenimento nei propri
    campioni biologici di qualsiasi sostanza vietata, metabolita o marker; ai fini
    dell'accertamento della violazione delle NSA non è necessario dimostrare il dolo, la colpa,
    la negligenza o l'uso consapevole da parte dell'Atleta;
  \item \textbf{2.1.2}: costituisce prova sufficiente di violazione uno dei seguenti casi:
    \begin{itemize}
      \item la presenza nel campione biologico A di una sostanza vietata o dei suoi metaboliti o
        marker, nel caso in cui l'Atleta rinunci all'analisi del campione biologico B e
        quest'ultimo non venga analizzato;
      \item la presenza nel campione biologico B di una sostanza vietata o dei suoi metaboliti o
        marker che confermi l'esito delle analisi effettuate sul campione biologico A: il
        campione biologico B dell'Atleta viene suddiviso in due flaconi e le analisi del secondo
        flacone confermano la presenza rinvenuta nel primo;
    \end{itemize}
  \item \textbf{2.1.3}: la mera presenza di un qualsiasi quantitativo di una sostanza vietata,
    dei suoi metaboliti o marker nel campione biologico dell'Atleta costituisce di per sé una
    violazione delle NSA, fatta eccezione per le sostanze per le quali la Lista delle sostanze e
    dei metodi proibiti indica specificamente un valore soglia;
  \item \textbf{2.1.4}: in deroga alla norma generale prevista dall'articolo 2.1, la Lista delle
    sostanze e dei metodi proibiti, ovvero gli Standard Internazionali, possono fissare alcuni
    criteri specifici per la valutazione delle sostanze vietate che possono essere prodotte per
    via endogena.
\end{itemize}

\subsection{2.2 -- Uso o tentato uso}

\textbf{Art.\ 2.2}: uso o tentato uso di una sostanza vietata o di un metodo proibito da parte
di un Atleta.
\begin{itemize}
  \item \textbf{2.2.1}: spetta a ogni Atleta accertarsi personalmente di non assumere alcuna
    sostanza vietata o di non utilizzare alcun metodo proibito; ai fini dell'accertamento della
    violazione delle NSA non sarà necessario dimostrare che vi sia dolo, colpa, negligenza o uso
    consapevole da parte dell'Atleta;
  \item \textbf{2.2.2}: il successo o il fallimento dell'uso o del tentato uso di una sostanza
    vietata o di un metodo proibito non costituiscono un elemento essenziale: è sufficiente,
    infatti, che la sostanza vietata o il metodo proibito siano stati usati o si sia tentato di
    usarli per integrare una violazione delle NSA.
\end{itemize}

\subsection{2.3 -- Elusione, rifiuto od omissione del prelievo}

\textbf{Art.\ 2.3}: eludere, rifiutarsi od omettere di sottoporsi al prelievo dei campioni
biologici; ovvero eludere il prelievo dei campioni biologici, oppure, senza giustificato motivo,
rifiutarsi di sottoporsi al prelievo previa notifica, in conformità alla normativa antidoping
applicabile.

\subsection{2.4 -- Mancata reperibilità}

\textbf{Art.\ 2.4}: mancata reperibilità (\textit{Whereabouts Failures}) --- ogni violazione
delle condizioni previste per gli Atleti che devono sottoporsi ai Controlli fuori competizione,
incluse la mancata presentazione di informazioni utili sulla reperibilità e la mancata esecuzione
di test in base a quanto previsto dal D-CI. Costituisce violazione ogni combinazione di 3 (tre)
mancati controlli e/o omesse comunicazioni, entro un periodo di 12 (dodici) mesi, da parte di un
Atleta inserito in un RTP.

\subsection{2.5 -- Manomissione}

\textbf{Art.\ 2.5}: manomissione o tentata manomissione in relazione a qualsiasi fase dei
Controlli antidoping --- condotta volta a minare la procedura di controllo antidoping ma che non
rientra nella definizione di pratiche vietate. La manomissione comprende, a titolo puramente
esemplificativo:
\begin{itemize}
  \item intralciare o tentare di intralciare intenzionalmente l'operato di un addetto al
    controllo antidoping;
  \item fornire informazioni fraudolente ad una Organizzazione Antidoping;
  \item intimidire o tentare di intimidire un potenziale testimone.
\end{itemize}

\subsection{2.6 -- Possesso}

\textbf{Art.\ 2.6}: possesso di sostanze vietate e ricorso a metodi proibiti.
\begin{itemize}
  \item \textbf{2.6.1}: possesso da parte di un Atleta, durante le competizioni, di qualsiasi
    sostanza vietata o il ricorso a qualsiasi metodo proibito, oppure possesso da parte di un
    Atleta, fuori competizione, di un metodo o di una sostanza espressamente vietati fuori
    competizione, a meno che l'Atleta possa dimostrare che il possesso sia dovuto ad un uso
    terapeutico consentito (artt.\ 14 e 15) o ad altro giustificato motivo;
  \item \textbf{2.6.2}: possesso da parte del Personale di supporto dell'Atleta, durante le
    competizioni, di qualsiasi sostanza vietata o di qualsiasi metodo proibito, oppure possesso
    da parte del Personale di supporto dell'Atleta, fuori competizione, di una sostanza o di un
    metodo espressamente vietati fuori competizione, in relazione a un Atleta, una competizione o
    un allenamento, a meno che il Personale possa dimostrare che il possesso sia dovuto ad un uso
    terapeutico consentito (artt.\ 14 e 15) o ad altro giustificato motivo.
\end{itemize}

\subsection{2.7 -- Traffico}

\textbf{Art.\ 2.7}: traffico illegale o tentato traffico illegale di sostanze vietate o metodi
proibiti.

\subsection{2.8 -- Somministrazione}

\textbf{Art.\ 2.8}: somministrazione o tentata somministrazione ad un Atleta durante le
competizioni, di una qualsiasi sostanza vietata o metodo proibito, oppure somministrazione o
tentata somministrazione ad un Atleta, fuori competizione, di una sostanza o di un metodo che
siano proibiti fuori competizione.

\subsection{2.9 -- Complicità}

\textbf{Art.\ 2.9}: fornire assistenza, incoraggiamento e aiuto, istigare, dissimulare o
assicurare ogni altro tipo di complicità intenzionale in riferimento a una qualsiasi violazione
o tentata violazione delle NSA o violazione dell'articolo 4.12.1 da parte di altra persona.

\subsection{2.10 -- Divieto di associazione}

\textbf{Art.\ 2.10}: l'Atleta o altra Persona soggetti all'autorità di una ADO, che si sia
avvalsa o abbia favorito la consulenza di Personale di supporto dell'Atleta:
\begin{itemize}
  \item \textbf{2.10.1}: soggetto all'autorità di una Organizzazione Antidoping che stia
    scontando un periodo di squalifica; ovvero
  \item \textbf{2.10.2}: non soggetto all'autorità di una Organizzazione Antidoping, che sia
    stato condannato o ritenuto colpevole solo nell'ambito di un procedimento penale,
    disciplinare o professionale per aver assunto una condotta che costituisca violazione del
    regolamento antidoping ove fossero stati applicati il Codice WADA e le NSA;
  \item \textbf{2.10.3}: funga da copertura o intermediario per un soggetto di cui agli artt.\
    2.10.1 o 2.10.2.
\end{itemize}

% ---------------------------------------------------------------------------
\section{Articolo 3 -- Altre violazioni delle Norme Sportive Antidoping}

Le seguenti voci costituiscono altre violazioni delle NSA:
\begin{itemize}
  \item \textbf{3.1}: qualsiasi violazione riferita alle fasi del controllo antidoping disposto
    dalla SVD di cui alla legge 376/2000;
  \item \textbf{3.2}: la mancata collaborazione da parte di qualunque soggetto per il rispetto
    delle NSA, ivi compresa l'omessa denuncia di circostanze rilevanti ai fini dell'accertamento
    di fatti di doping;
  \item \textbf{3.3}: la condotta offensiva nei confronti del DCO e/o del Personale addetto al
    controllo antidoping, la quale non sia configurabile come violazione dell'articolo 2.5 del
    CSA.
\end{itemize}

% ---------------------------------------------------------------------------
\section{Il Disciplinare per le Esenzioni ai Fini Terapeutici (D-EFT)}

\textbf{Disciplinare per le Esenzioni ai Fini Terapeutici} (\textbf{D-EFT}): attuativo
dell'\textit{International Standard for Therapeutic Use Exemption} WADA.

\textbf{Esenzione a fini terapeutici} (\textit{Therapeutic Use Exemption}, \textbf{TUE}): un
Atleta che, per motivi di salute, deve ricorrere all'assunzione e/o somministrazione di farmaci
contenenti principi attivi compresi nella Lista, deve presentare una domanda di TUE.
\begin{itemize}
  \item gli atleti di livello \textbf{Nazionale} devono presentare la domanda di TUE al
    \textbf{CEFT};
  \item gli atleti di livello \textbf{Internazionale} devono presentare la domanda di TUE al
    \textbf{TUEC} della Federazione Internazionale (FI) di appartenenza.
\end{itemize}

\subsection{Criteri per la concessione di una TUE}

\begin{itemize}
  \item \textbf{art.\ 4.1 a}: l'Atleta potrebbe subire un danno se la terapia farmacologica
    venisse sospesa;
  \item \textbf{art.\ 4.1 b}: l'uso della sostanza non deve produrre un incremento della
    performance;
  \item \textbf{art.\ 4.1 c}: non è possibile ricorrere ad una terapia farmacologica con farmaci
    non proibiti;
  \item \textbf{art.\ 4.1 d}: la necessità di utilizzare la sostanza non può essere conseguenza
    di un precedente utilizzo, senza TUE, di sostanze proibite.
\end{itemize}

\rubrica{Durata ed estinzione della TUE}
\begin{itemize}
  \item \textbf{1.2}: ciascuna TUE ha una precisa durata, così come decisa dal CEFT, al termine
    della quale cessa automaticamente di avere efficacia; se l'Atleta necessita di proseguire
    l'utilizzo della sostanza o del metodo proibiti dopo la scadenza della TUE, deve presentare,
    prima della scadenza, una nuova domanda di TUE;
  \item \textbf{1.3}: se una TUE è scaduta o è stata revocata e la sostanza proibita soggetta
    alla TUE è ancora presente nell'organismo dell'Atleta, la \textbf{Procura Nazionale
    Antidoping} (\textbf{PNA}), a seguito di un riscontro di Esito Avverso, interpella il CEFT,
    che valuta se il referto è compatibile con la scadenza o la revoca della TUE.
\end{itemize}

\subsection{Procedura per la presentazione di una domanda di TUE}

La domanda di TUE prevede la trasmissione al CEFT, a mezzo raccomandata con ricevuta di ritorno
anticipata via fax, ovvero a mezzo posta elettronica. Deve essere presentata al CEFT
\textbf{prima} di cominciare una terapia.
\begin{itemize}
  \item per le sostanze proibite in e fuori competizione, la domanda deve essere presentata
    appena formulata la diagnosi che preveda l'utilizzo di sostanze o metodi proibiti;
  \item nei casi di condizione clinica acuta e/o non procrastinabile, la domanda deve essere
    trasmessa tempestivamente \textbf{dopo} l'inizio e/o l'effettuazione (1 sola
    somministrazione) della terapia (\textbf{TUE retroattiva}).
\end{itemize}

\rubrica{Documentazione richiesta}
\begin{itemize}
  \item Modulo TUE F49 (\textit{Therapeutic Use Exemption Application});
  \item scheda per il medico curante/specialista, Modulo F51;
  \item anamnesi, storia clinica medica e documentazione comprovante la diagnosi, comprensiva dei
    risultati degli accertamenti specifici della patologia in essere e della diagnostica per
    immagini, e certificazione del medico specialista nella patologia di cui trattasi, che
    attesti sia l'assenza di eventuali controindicazioni, anche temporanee, alla pratica
    dell'attività sportiva agonistica, sia la necessità dell'utilizzo della sostanza o del metodo
    proibiti nella cura dell'Atleta, motivando le ragioni per cui non è possibile utilizzare un
    altro farmaco consentito;
  \item certificato medico di idoneità all'attività agonistica, ove previsto dalla normativa
    sportiva della disciplina di appartenenza o da norme di legge, e/o, per gli atleti
    professionisti di cui alla legge 91/1981, scheda sanitaria aggiornata con riferimento alla
    patologia per cui si richiede la TUE;
  \item certificato medico per l'attività sportiva non agonistica, ove previsto per la disciplina
    praticata.
\end{itemize}

La modulistica deve essere compilata in ogni sua parte, con redazione dattilografica o in
stampatello (``capital letter''), specificando:
\begin{itemize}
  \item Federazione Sportiva Nazionale (\textbf{FSN}) / Disciplina Sportiva Associata
    (\textbf{DSA}) / Ente di Promozione Sportiva (\textbf{EPS}) di appartenenza e la disciplina
    sportiva praticata dall'Atleta;
  \item diagnosi;
  \item principi attivi contenuti in medicinali registrati (\textit{generic name}), via di
    somministrazione (\textit{route}), dosaggio (\textit{dose}), posologia (\textit{frequency});
  \item durata di somministrazione della sostanza o dell'applicazione del metodo normalmente
    vietati per cui si richiede l'esenzione, specificando la data di inizio e la data di fine
    dell'intervento farmacologico.
\end{itemize}
La modulistica illeggibile o ritenuta incompleta non sarà esaminata e verrà restituita
all'interessato.

\subsection{Il certificato di idoneità all'attività sportiva agonistica}

Ai sensi dell'articolo 1 comma 4 della legge 376/2000, nonché delle norme per la tutela
sanitaria dell'attività sportiva agonistica contenute nei regolamenti sanitari sportivi, è cura
del medico che rilascia il certificato di idoneità informare l'Atleta in ordine agli obblighi di
conservazione di tutta la propria documentazione medica, per eventuali richieste delle Autorità
sportive.

Le esenzioni concesse dal CEFT sono comunque subordinate al rilascio, alla vigenza ovvero alla
riemissione su richiesta del CEFT stesso, del certificato di idoneità sportiva agonistica, e
comportano l'aggiornamento della scheda sanitaria per gli atleti professionisti, a norma
dell'articolo 7 della legge 23 marzo 1981, n.\,91.

\subsection{Decisione del CEFT e riconoscimento internazionale}

L'Atleta deve attendere la decisione del CEFT prima di poter fare uso della sostanza o del
metodo per cui ha richiesto la TUE.

\rubrica{Procedura e criteri di riconoscimento internazionale di una TUE}
\begin{itemize}
  \item un Atleta di livello Internazionale deve presentare la domanda di TUE alla Federazione
    Internazionale di appartenenza;
  \item laddove un Atleta di livello Internazionale abbia già una TUE concessa dal CEFT, la
    rispettiva Federazione Internazionale di appartenenza deve riconoscerne la validità.
\end{itemize}
